
\chapter*{EXPOSÉ XVII\\[0.4em]
UNIPOTENT ALGEBRAIC GROUPS. EXTENSIONS BETWEEN UNIPOTENT GROUPS AND GROUPS OF MULTIPLICATIVE TYPE}
\addcontentsline{toc}{chapter}{Exposé XVII. Unipotent Algebraic Groups. Extensions}

\begin{center}
\textit{by M. Raynaud}%
\begingroup
\renewcommand{\thefootnote}{*}%
\footnote{Cf. the note on page 1 of Exposé XV.}%
\endgroup
\end{center}

\section*{0.\ Some notation%
\begingroup
\renewcommand{\thefootnote}{0}%
\endgroup}
\addcontentsline{toc}{section}{0.\ Some notation}
\SGARefTarget{sga3:xvii:section:0}

In this chapter we shall be concerned chiefly with algebraic groups defined over a
field \(k\).  The number \(p\geq 0\) will always denote the characteristic of \(k\);
\(\mathbf F_p\) will denote the prime field with \(p\) elements when \(p>0\);
\(\bar{k}\) will denote an algebraically closed extension of \(k\); and \(q\) will
denote a prime number different from \(p\).

For every \(S\)-prescheme, \((\mathbb G_a)_S\) (respectively
\((\mathbb G_m)_S\)) denotes the additive group (respectively the multiplicative
group) over \(S\) (cf.\ \SGARefLink{sga3:i:subsection:4-3}{Exposé I, 4.3}).  For every integer \(n>0\),
\((\boldsymbol\mu_n)_S\) (respectively \((\mathbb Z/n\mathbb Z)_S\)) denotes the
group of \(n\)-th roots of unity (\SGARefLink{sga3:i:paragraph:4-4-4}{Exposé I, 4.4.4}) (respectively the constant group
over \(S\) associated with the abstract group \(\mathbb Z/n\mathbb Z\)
(\SGARefLink{sga3:i:subsection:4-1}{Exposé I, 4.1})).  The group \((\mathbb Z/n\mathbb Z)_S\) is finite and étale
over \(S\); the group \((\boldsymbol\mu_n)_S\) is flat and finite over \(S\), and
is étale over \(S\) if and only if \(n\) is invertible on \(S\)
(\SGARefLink{sga3:viii:proposition:2:1}{Exposé VIII, 2.1}).

Suppose that \(S\) is a prescheme of characteristic \(p>0\).  For every integer
\(n>0\) and every \(S\)-group prescheme \(G\), we denote by \(F^n(G)\) the
radicial \(S\)-subgroup prescheme of \(G\) equal to the kernel of the \(n\)-th
iterate of the relative Frobenius morphism of \(G\) (Exposé VII\textsubscript{A}).
In particular, if \(G=(\mathbb G_a)_S\), we set
\[
  F^n(G)=(\boldsymbol\alpha_{p^n})_S .
\]
This is a radicial \(S\)-group, flat and finite over \(S\), representing the
following functor: for every \(S\)-prescheme \(S'\),
\((\boldsymbol\alpha_{p^n})_S(S')\) is the set of
\(x'\in\Gamma(S',\mathcal O_{S'})\) such that \(x'^{p^n}=0\).

The already defined group \((\boldsymbol\mu_{p^n})_S\) is canonically isomorphic
to \(F^n((\mathbb G_m)_S)\).

When there is no ambiguity about the base scheme \(S\), we shall simply write
\(\mathbb G_a\), \(\mathbb G_m\), \(\boldsymbol\alpha_{p^n}\), and so on, in
place of \((\mathbb G_a)_S\), \((\mathbb G_m)_S\),
\((\boldsymbol\alpha_{p^n})_S\), and so on.

If \(G\) is a commutative \(S\)-group prescheme, then for every integer \(n>0\),
\({}_nG\) denotes the subgroup prescheme of \(G\) equal to the kernel of
raising to the \(n\)-th power in \(G\).

For the reader's convenience, we have collected in an appendix some properties
of algebraic groups proved in Exposés VI and VII, together with elementary
properties of Hochschild cohomology that will be useful in this chapter.
